% Luc Destombe --- licence CC BY-NC-SA 4.0
% https://creativecommons.org/licenses/by-nc-sa/4.0/deed.fr
% Fichier autonome : compiler avec pdflatex, deux fois (signets, nombre de pages).
\documentclass[11pt,a4paper]{article}
\makeatletter
\newif\iflycee@unecolonne
\lycee@unecolonnetrue
\newif\iflycee@cm
\newif\iflycee@palatino

\RequirePackage[utf8]{inputenc}
\RequirePackage[T1]{fontenc}
\RequirePackage[french]{babel}
\RequirePackage{etoolbox}

\newcommand{\ShorthandsOff}{\@ifundefined{shorthandoff}{}{\shorthandoff{;:!?}}}
\newcommand{\ShorthandsOn}{\@ifundefined{shorthandon}{}{\shorthandon{;:!?}}}
\ShorthandsOff
\AtBeginDocument{\ShorthandsOn}
\AtBeginEnvironment{tikzpicture}{\ShorthandsOff}

\RequirePackage{geometry}
\RequirePackage{fancyhdr}
\RequirePackage{lastpage}
\RequirePackage{multicol}
\RequirePackage{array}
\RequirePackage{enumitem}
\RequirePackage{xcolor}
\RequirePackage{needspace}

\RequirePackage{amsmath,amssymb}
\iflycee@cm\else\RequirePackage{fourier}\fi
\DeclareMathAlphabet{\mathcal}{OMS}{cmsy}{m}{n}
\renewcommand{\ttdefault}{lmtt}
\begingroup\catcode`\_=\active \catcode`\^=\active
\global\let_\sb \global\let^\sp
\endgroup
\AtBeginDocument{\catcode`\_=12 \mathcode`\_="8000
  \catcode`\^=12 \mathcode`\^="8000 }
\newlength{\retraitcalculs}
\setlength{\retraitcalculs}{2em}

\RequirePackage{tikz}
\usetikzlibrary{arrows.meta,calc,babel,shapes.misc}
\RequirePackage[most]{tcolorbox}
\tcbuselibrary{skins,breakable}

\definecolor{coulDef}{HTML}{1F4E79}
\definecolor{coulProp}{HTML}{7030A0}
\definecolor{coulRem}{HTML}{C55A11}
\definecolor{coulEx}{HTML}{2E7D32}
\definecolor{coulExo}{HTML}{B03A2E}
\definecolor{coulPart}{HTML}{1F4E79}
\definecolor{coulMeth}{HTML}{1B6E4F}
\definecolor{coulCourbe}{HTML}{0B5ED7}
\definecolor{coulCourbeB}{HTML}{C00000}
\definecolor{coulCourbeC}{HTML}{2E7D32}
\definecolor{coulGrille}{HTML}{8C8C8C}
\definecolor{coulRep}{HTML}{1B6E4F}
\definecolor{coulBar}{HTML}{2E7D32}

\newcommand{\R}{\mathbb{R}}
\newcommand{\Rs}{\mathbb{R}^{*}}
\newcommand{\N}{\mathbb{N}}
\newcommand{\Ns}{\mathbb{N}^{*}}
\newcommand{\Z}{\mathbb{Z}}
\newcommand{\Q}{\mathbb{Q}}
\renewcommand{\le}{\leqslant}
\renewcommand{\ge}{\geqslant}
\renewcommand{\approx}{\simeq}
\newcommand{\vect}[1]{\overrightarrow{#1}}
\newcommand{\norme}[1]{\left\lVert #1 \right\rVert}
\newcommand{\intervalle}[2]{\left[#1\,;#2\right]}
\RequirePackage{cancel}
\renewcommand{\CancelColor}{\color{coulExo}}
\newcommand{\fois}[1]{\times{\color{coulExo}#1}}
\newcommand{\barre}[1]{\cancel{#1}}
\newlength{\lycee@eqh}\newlength{\lycee@eqd}
\newsavebox{\lycee@eqa}\newsavebox{\lycee@eqb}\newsavebox{\lycee@eqc}
\newcommand{\lycee@eqcol}[1]{\hbox{$\renewcommand{\arraystretch}{1.35}%
  \begin{array}[t]{@{}>{\displaystyle}r@{}>{\displaystyle{}}l@{}}#1\end{array}$}}
\newcommand{\lycee@eqmesure}[1]{%
  \ifdim\ht#1>\lycee@eqh\setlength{\lycee@eqh}{\ht#1}\fi
  \ifdim\dp#1>\lycee@eqd\setlength{\lycee@eqd}{\dp#1}\fi}
\newcommand{\lycee@eqtrait}{\hspace{1.2em}ou\hspace{1.2em}}
\newcommand{\lycee@eqiff}{$\iff$\hspace{0.8em}}
\newcommand{\eqproduit}[2]{\par\smallskip\noindent C'est une équation produit :
  $A\times B=0 \iff A=0$ ou $B=0$.\par\nopagebreak\smallskip
  \begingroup\setlength{\lycee@eqh}{0pt}\setlength{\lycee@eqd}{0pt}%
  \setbox\lycee@eqa=\lycee@eqcol{#1}\setbox\lycee@eqb=\lycee@eqcol{#2}%
  \lycee@eqmesure\lycee@eqa\lycee@eqmesure\lycee@eqb%
  \noindent\hspace*{\retraitcalculs}\lycee@eqiff\usebox\lycee@eqa\lycee@eqtrait\usebox\lycee@eqb\endgroup\par\smallskip}
\newcommand{\eqproduittrois}[3]{\par\smallskip\noindent C'est une équation produit :
  $A\times B\times C=0 \iff A=0$ ou $B=0$ ou $C=0$.\par\nopagebreak\smallskip
  \begingroup\setlength{\lycee@eqh}{0pt}\setlength{\lycee@eqd}{0pt}%
  \setbox\lycee@eqa=\lycee@eqcol{#1}\setbox\lycee@eqb=\lycee@eqcol{#2}\setbox\lycee@eqc=\lycee@eqcol{#3}%
  \lycee@eqmesure\lycee@eqa\lycee@eqmesure\lycee@eqb\lycee@eqmesure\lycee@eqc%
  \noindent\hspace*{\retraitcalculs}\lycee@eqiff\usebox\lycee@eqa\lycee@eqtrait\usebox\lycee@eqb\lycee@eqtrait\usebox\lycee@eqc\endgroup
  \par\smallskip}

\tikzset{grille/.style={coulGrille,line width=0.4pt}}
\newcommand{\quadrillage}[4]{\draw[grille,step=1] (#1,#2) grid (#3,#4);}
\tikzset{
  croix/.style={cross out,draw,line width=0.6pt,inner sep=0pt,minimum size=4pt},
  croixdroite/.style={croix,rotate=45,line width=0.8pt},
}
\newcommand{\pt}[3]{\path (#1) node[croix]{} node[#2,font=\small,inner sep=2.5pt]{$#3$};}

\newcommand{\lycee@repere}[4]{%
  \draw[grille,step=1] (#1,#2) grid (#3,#4);
  \draw[-{Stealth[length=2mm]},thick] (#1,0)--({#3+0.4},0) node[below left=-1pt]{$x$};
  \draw[-{Stealth[length=2mm]},thick] (0,#2)--(0,{#4+0.4}) node[below left=-1pt]{$y$};
  \node[below left,font=\scriptsize,inner sep=1.5pt] at (0,0) {$0$};}
\newcommand{\lycee@gradx}[1]{\foreach \k/\l in {#1}{%
  \draw (\k,0.07)--(\k,-0.07) node[below,font=\scriptsize,inner sep=1.5pt]{$\l$};}}
\newcommand{\lycee@grady}[1]{\foreach \k/\l in {#1}{%
  \draw (0.07,\k)--(-0.07,\k) node[left,font=\scriptsize,inner sep=1.5pt]{$\l$};}}
\newcommand{\lycee@intff}[2]{\left[#1\,;#2\right]}
\newcommand{\lycee@intfo}[2]{\left[#1\,;#2\right[}
\newcommand{\lycee@intof}[2]{\left]#1\,;#2\right]}
\newcommand{\lycee@intoo}[2]{\left]#1\,;#2\right[}
\newcommand{\lycee@ens}[1]{\left\{#1\right\}}
\AtBeginDocument{%
  \providecommand{\repere}{\lycee@repere}%
  \providecommand{\gradx}{\lycee@gradx}%
  \providecommand{\grady}{\lycee@grady}%
  \providecommand{\Cf}{\mathcal{C}\sb{f}}%
  \providecommand{\Cg}{\mathcal{C}\sb{g}}%
  \providecommand{\intff}{\lycee@intff}%
  \providecommand{\intfo}{\lycee@intfo}%
  \providecommand{\intof}{\lycee@intof}%
  \providecommand{\intoo}{\lycee@intoo}%
  \providecommand{\ens}{\lycee@ens}}

\newlist{questions}{enumerate}{3}
\setlist[questions,1]{label=\arabic*\textdegree),leftmargin=*,itemsep=2pt,topsep=3pt}
\setlist[questions,2]{label=\alph*),leftmargin=*,itemsep=1pt,topsep=2pt}
\setlist[questions,3]{label=\roman*),leftmargin=*,itemsep=1pt,topsep=1pt}
\setlist[itemize]{label=\textbullet,leftmargin=*,itemsep=2pt,topsep=3pt}

\newcommand{\besoinplace}[1]{\par\penalty\@M\begingroup
  \dimen@=#1\relax\dimen@ii=\pagegoal\advance\dimen@ii-\pagetotal
  \ifdim\dimen@>\dimen@ii\ifdim\dimen@ii>\z@\vfil\fi\break\fi\endgroup}

\newcommand{\lycee@pied}{\thepage/\pageref{LastPage}}
\newcommand{\entete}[2]{%
  \pagestyle{fancy}\fancyhf{}%
  \fancyhead[L]{\footnotesize\color{coulPart}\textbf{#1}}%
  \fancyhead[R]{\footnotesize\color{coulPart}\textbf{#2}}%
  \fancyfoot[C]{\footnotesize\lycee@pied}%
  \renewcommand{\headrulewidth}{0.4pt}%
  \renewcommand{\headrule}{\hbox to\headwidth{%
    \color{coulPart!40}\leaders\hrule height \headrulewidth\hfill}}}

\RequirePackage{ccicons}
\newcommand{\lycee@licence}{{\scriptsize\color{gray}%
  \copyright\ Luc Destombe --- \ccbyncsa\ CC BY-NC-SA 4.0}}
\newif\iflycee@licence \lycee@licencetrue
\newcommand{\sanslicence}{\lycee@licencefalse}
\AtBeginDocument{\iflycee@licence\AddToHookNext{shipout/foreground}{%
  \put(0,0){\rlap{\kern\dimexpr1in+\hoffset+\oddsidemargin+\textwidth\relax
    \llap{\raisebox{-\dimexpr1in+\voffset+\topmargin+\headheight+\headsep
      +\textheight+\footskip\relax}{\lycee@licence}}}}}\fi}

\newcommand{\lycee@site}{faire-des-maths.fr}
\newcommand{\lycee@version}{{\scriptsize\color{gray}\lycee@site
  \ifdefined\lyceeversion\ --- Version~\lyceeversion\fi}}
\AtBeginDocument{\iflycee@licence\AddToHookNext{shipout/foreground}{%
  \put(0,0){\rlap{\kern\dimexpr1in+\hoffset+\oddsidemargin\relax
    \raisebox{-\dimexpr1in+\voffset+\topmargin+\headheight+\headsep
      +\textheight+\footskip\relax}{\lycee@version}}}}\fi}

\newcommand{\formulesserrees}{%
  \setlength{\abovedisplayskip}{3pt}\setlength{\belowdisplayskip}{3pt}%
  \setlength{\abovedisplayshortskip}{2pt}\setlength{\belowdisplayshortskip}{3pt}}

\setlength{\parindent}{0pt}

\RequirePackage[implicit=false,hidelinks,pdfencoding=auto,psdextra,bookmarksopen,
  bookmarksopenlevel=1,pdfcreator={},pdfproducer={}]{hyperref}
\RequirePackage{bookmark}
\hypersetup{pdfauthor={Luc Destombe},pdfkeywords={CC BY-NC-SA 4.0}}
\newcounter{lycee@signet}
\newcommand{\signet}[3][]{\stepcounter{lycee@signet}%
  \edef\lycee@dest{\if\relax\detokenize{#1}\relax
    signet.\arabic{lycee@signet}\else#1\fi}%
  \hypertarget{\lycee@dest}{}\bookmark[level=#2,dest=\lycee@dest]{#3}}
\pdfstringdefDisableCommands{%
  \def\R{R}\def\N{N}\def\Z{Z}\def\Q{Q}\def\vect#1{#1}\def\Cf{Cf}\def\Cg{Cg}\def\fois#1{×#1}%
  \def\barre#1{#1}\def\intervalle#1#2{[#1;#2]}\def\intff#1#2{[#1;#2]}%
  \def\textsuperscript#1{#1}\def\dfrac#1#2{#1/#2}\def\frac#1#2{#1/#2}%
  \def\vec#1{#1⃗}}%
\RequirePackage[official]{eurosym}

\tcbset{exercice corrige/.style={
  enhanced,breakable,before upper=\raggedright,
  colback=white,colframe=coulExo,
  boxrule=0.8pt,arc=2pt,
  left=6pt,right=6pt,top=8pt,bottom=5pt,
  before skip=9pt,after skip=4pt,
  coltitle=white,fonttitle=\bfseries\small,
  attach boxed title to top left={xshift=8pt,yshift=-\tcboxedtitleheight/2},
  boxed title style={colback=coulExo,arc=2pt,boxrule=0pt}}}
\newcounter{exo}
\newtcolorbox{exercice}[1][]{exercice corrige,
  phantom={\signet[exercice-\arabic{exo}]{2}{Exercice \arabic{exo}}},
  title={Exercice \theexo\ifx\relax#1\relax\else{} --- #1\fi}}
\newenvironment{exo}[1][\relax]{\refstepcounter{exo}\begin{exercice}[#1]}{\end{exercice}}

\geometry{top=1.8cm,bottom=1cm,left=1cm,right=1cm,headheight=14pt,footskip=0.6cm}
\renewcommand{\lycee@pied}{\thepage}

\renewcommand{\theexo}{\ifnum\value{exo}<10 0\fi\arabic{exo}}
\tcbset{exercice corrige/.style={
  enhanced,breakable,before upper=\raggedright,
  colback=white,colframe=coulExo,
  boxrule=0.9pt,arc=2pt,
  left=8pt,right=8pt,top=8pt,bottom=6pt,
  before skip=8pt,after skip=6pt,
  coltitle=white,fonttitle=\bfseries,
  attach boxed title to top left={xshift=8pt,yshift=-\tcboxedtitleheight/2},
  boxed title style={colback=coulExo,arc=2pt,boxrule=0pt}}}

\newcommand{\entetefiche}[2]{%
  \begin{center}
    \begin{tcolorbox}[enhanced,width=0.92\linewidth,colback=white,colframe=coulPart,
        boxrule=1.4pt,arc=3pt,halign=center,top=6pt,bottom=6pt]
      {\large\bfseries\color{coulPart} CORRIGÉ --- FICHE D'EXERCICES}\\[3pt]
      {\normalsize\bfseries #1}\\[2pt]
      {\normalsize #2}
    \end{tcolorbox}
  \end{center}}

\newcounter{partie}
\newcommand{\partie}[1]{%
  \stepcounter{partie}%
  \besoinplace{8\baselineskip}\vspace{4pt}%
  \begin{tcolorbox}[enhanced,colback=coulPart!8,colframe=coulPart,boxrule=0pt,
      leftrule=3pt,arc=0pt,left=6pt,right=4pt,top=3pt,bottom=3pt,
      before skip=6pt,after skip=2pt]
    \bfseries\color{coulPart}\Roman{partie} --- #1\signet{1}{\Roman{partie} --- #1}
  \end{tcolorbox}}
\newcommand{\partiesansnum}[1]{%
  \besoinplace{8\baselineskip}\vspace{4pt}%
  \begin{tcolorbox}[enhanced,colback=coulPart!8,colframe=coulPart,boxrule=0pt,
      leftrule=3pt,arc=0pt,left=6pt,right=4pt,top=3pt,bottom=3pt,
      before skip=6pt,after skip=2pt]
    \bfseries\color{coulPart}#1\signet{1}{#1}
  \end{tcolorbox}}

\setlist[questions,1]{label=\arabic*\textdegree),leftmargin=*,itemsep=3pt,topsep=3pt}
\setlist[questions,2]{label=\alph*),leftmargin=*,itemsep=2pt,topsep=2pt}
\setlist[questions,3]{label=\roman*),leftmargin=*,itemsep=1pt,topsep=1pt}
\newlist{expr}{enumerate}{1}
\setlist[expr]{label={\Alph*\ $=$},leftmargin=*,itemsep=2pt,topsep=3pt}

\newcommand{\res}[1]{{\color{coulRep}#1}}
\newcommand{\rem}[1]{\par\smallskip{\small\itshape\color{black!60}$\rightarrow$ #1}}
\newenvironment{calculs}{\par\smallskip\noindent$\displaystyle\begin{aligned}[t]}%
  {\end{aligned}$\par\smallskip}
\newcommand{\justif}[1]{\quad\text{\small\itshape\color{black!60}#1}}
\newcommand{\cas}[1]{\par\smallskip\noindent\textbf{\color{coulExo}#1}\vadjust{\nobreak}}
\newcommand{\explique}[2]{\par\smallskip\noindent
  \begin{minipage}[t]{0.6\linewidth}\raggedright #1\end{minipage}\hfill
  \begin{minipage}[t]{0.37\linewidth}\raggedright\leavevmode
    {\small\itshape\color{black!60}#2}\end{minipage}\par}

\setlength{\parskip}{2pt}

\makeatother
\entete{Chapitre 10 --- Probabilités conditionnelles --- corrigé}{1\textsuperscript{re} STI2D}
\let\casligne\cas
\renewcommand{\cas}[1]{\casligne{#1}\par\nopagebreak}
\let\partieligne\partie
\renewcommand{\partie}[1]{\besoinplace{16\baselineskip}\partieligne{#1}}
\BeforeBeginEnvironment{exo}{\besoinplace{10\baselineskip}}
\setlength{\parindent}{0pt}
\setlength{\jot}{5pt}

\newcommand{\venn}[6]{%
  \draw (-1.75,-1.35) rectangle (2.75,1.35);
  \node[font=\footnotesize] at (-1.5,1.1) {$\Omega$};
  \draw[coulCourbe,thick] (0,0) circle (1.1);
  \draw[coulCourbeB,thick] (1,0) circle (1.1);
  \node[coulCourbe] at (-0.95,0.95) {$#1$};
  \node[coulCourbeB] at (1.95,0.95) {$#2$};
  \node at (-0.6,0) {#3};
  \node at (0.5,0) {#4};
  \node at (1.6,0) {#5};
  \node at (2.15,-1.05) {#6};}
\newcommand{\vennliste}[6]{%
  \draw (-1.75,-1.95) rectangle (2.75,1.35);
  \node[font=\footnotesize] at (-1.5,1.1) {$\Omega$};
  \draw[coulCourbe,thick] (0,0) circle (1.1);
  \draw[coulCourbeB,thick] (1,0) circle (1.1);
  \node[coulCourbe] at (-0.95,0.95) {$#1$};
  \node[coulCourbeB] at (1.95,0.95) {$#2$};
  \node[font=\footnotesize,align=center] at (-0.55,0) {#3};
  \node[font=\footnotesize,align=center] at (0.5,0) {#4};
  \node[font=\footnotesize,align=center] at (1.55,0) {#5};
  \node[font=\footnotesize] at (0.5,-1.6) {#6};}
\tikzset{evt/.style={inner sep=2pt},
  proba/.style={font=\footnotesize,fill=white,inner sep=1pt}}
\newcommand{\arbre}[8]{%
  \node (O) at (0,0) {};
  \node[evt] (A) at (2.6,1.4) {$#1$};
  \node[evt] (nA) at (2.6,-1.4) {$\overline{#1}$};
  \node[evt] (AB) at (5.2,2.1) {$#2$};
  \node[evt] (AnB) at (5.2,0.7) {$\overline{#2}$};
  \node[evt] (nAB) at (5.2,-0.7) {$#2$};
  \node[evt] (nAnB) at (5.2,-2.1) {$\overline{#2}$};
  \draw (O.center) -- node[proba]{$#3$} (A);
  \draw (O.center) -- node[proba]{$#4$} (nA);
  \draw (A) -- node[proba]{$#5$} (AB);
  \draw (A) -- node[proba]{$#6$} (AnB);
  \draw (nA) -- node[proba]{$#7$} (nAB);
  \draw (nA) -- node[proba]{$#8$} (nAnB);}
\newcommand{\feuille}[2]{\node[right,font=\footnotesize] at (#1.east) {#2};}
\newcommand{\feuilles}[2]{\feuille{AB}{$#1\cap #2$}\feuille{AnB}{$#1\cap\overline{#2}$}%
  \feuille{nAB}{$\overline{#1}\cap #2$}\feuille{nAnB}{$\overline{#1}\cap\overline{#2}$}}

\begin{document}

\entetefiche{Chapitre 10 --- Probabilités conditionnelles}{Première STI2D}

\partie{Vocabulaire des événements}

\begin{exo}[Décrire des événements]
\cas{1)}
\explique{$\Omega=\res{\ens{1\,;2\,;3\,;\dots\,;20}}$ : $20$ issues.}{tous les numéros des
jetons}
\explique{$A=\res{\ens{4\,;8\,;12\,;16\,;20}}$}{les multiples de $4$ jusqu'à $20$}
\explique{$B=\res{\ens{15\,;16\,;17\,;18\,;19\,;20}}$}{$15$ compris : \og supérieur ou
égal\fg{}}
\cas{2)}
\explique{$A\cap B$ : \og obtenir un multiple de $4$ \textbf{et} un nombre supérieur ou égal
à $15$\fg{} ; $A\cap B=\res{\ens{16\,;20}}$.}{les issues communes à $A$ et à $B$}
\explique{$A\cup B$ : \og obtenir un multiple de $4$ \textbf{ou} un nombre supérieur ou égal
à $15$\fg{} ;\\ $A\cup B=\res{\ens{4\,;8\,;12\,;15\,;16\,;17\,;18\,;19\,;20}}$.}{les issues
de $A$, puis celles de $B$ qui manquent ; $16$ et $20$ ne s'écrivent qu'une fois}
\cas{3)}
\explique{$\overline{B}=\res{\ens{1\,;2\,;3\,;\dots\,;14}}$ : \og obtenir un nombre
\res{inférieur ou égal à $14$}\fg{}.}{les issues de $\Omega$ qui ne sont pas dans $B$}
\cas{4)}
\explique{$C=\ens{1\,;3\,;5\,;\dots\,;19}$ et $A\cap C=\varnothing$, donc \res{$A$ et $C$
sont incompatibles}.}{un multiple de $4$ est pair, jamais impair}
\explique{$B\cap C=\ens{15\,;17\,;19}$, donc \res{$B$ et $C$ ne sont pas incompatibles}.}{si
on obtient $15$, $B$ et $C$ sont réalisés en même temps}
\end{exo}

\begin{exo}[Écrire l'univers]
\cas{a)}
\explique{$\Omega=\res{\ens{\text{P}\,;\text{F}}}$}{deux issues}
\cas{b)}
\explique{$\Omega=\res{\ens{\text{R}\,;\text{V}\,;\text{B}\,;\text{J}}}$}{une issue par
secteur}
\cas{c)}
\explique{$\Omega=\res{\ens{11\,;12\,;21\,;22}}$}{$2$ choix pour le premier chiffre, puis $2$
pour le second}
\cas{d)}
\explique{$\Omega=\res{\ens{1\,;2\,;6}}$}{on note le nombre lu : trois résultats
possibles, même si le dé a six faces}
\cas{e)}
\explique{$\Omega=\res{\ens{\text{L}\,;\text{M}\,;\text{J}\,;\text{V}\,;\text{S}\,;\text{D}}}$
: $6$ issues.}{mardi et mercredi commencent tous deux par M : on ne l'écrit qu'une fois}
\rem{Dans d) et e), une issue s'écrit une seule fois dans l'univers, même quand elle peut
se produire de plusieurs façons.}
\end{exo}

\begin{exo}[En phrase complète]
\cas{1)}
\explique{$A\cap B$ : \og \res{l'élève est en première \textbf{et} fait partie de
l'association sportive}\fg{}.}{les deux à la fois}
\explique{$A\cup B$ : \og \res{l'élève est en première \textbf{ou} fait partie de
l'association sportive}\fg{} (ou les deux).}{au moins un des deux}
\explique{$\overline{A}$ : \og \res{l'élève n'est pas en première}\fg{}.}{il est en
seconde ou en terminale}
\explique{$\overline{A}\cap B$ : \og \res{l'élève fait partie de l'association sportive,
mais n'est pas en première}\fg{}.}{$B$ \textbf{et} le contraire de $A$}
\cas{2)}
\explique{L'élève fait partie de l'association sportive : \res{$B$ est réalisé}, donc
\res{$A\cup B$ aussi}.}{pour $A\cup B$, il suffit qu'un des deux soit réalisé}
\explique{Il est en seconde, donc pas en première : \res{$\overline{A}$ est réalisé}, mais
\res{$A$ et $A\cap B$ ne le sont pas}.}{pour $A\cap B$, il faudrait qu'il soit aussi en
première}
\end{exo}

\begin{exo}[Des mots dans une urne]
\cas{1)}
\explique{$A=\res{\ens{\text{tram}\,;\text{car}\,;\text{train}\,;\text{métro}\,;\text{ferry}}}$}{les
mots qui contiennent un r}
\explique{$B=\res{\ens{\text{vélo}\,;\text{tram}\,;\text{taxi}}}$}{bus et car ont trois
lettres ; train, métro et ferry en ont cinq}
\cas{2)}
\explique{$A\cap B=\res{\ens{\text{tram}}}$}{le seul mot de quatre lettres qui contient un
r}
\explique{$A\cup B=\res{\ens{\text{vélo}\,;\text{tram}\,;\text{car}\,;\text{train}\,;\text{métro}\,;\text{taxi}\,;\text{ferry}}}$}{tous
les mots sauf bus}
\cas{3)}
\explique{$\overline{A}=\res{\ens{\text{vélo}\,;\text{bus}\,;\text{taxi}}}$ : \og \res{le mot
ne contient pas la lettre r}\fg{}.}{les mots de $\Omega$ qui ne sont pas dans $A$}
\cas{4)}
\explique{$\overline{A}\cap\overline{B}=\res{\ens{\text{bus}}}$ : \og \res{le mot ne contient
pas la lettre r et n'a pas quatre lettres}\fg{}.}{parmi vélo, bus et taxi, seul bus n'a pas
quatre lettres}
\end{exo}

\partie{Diagramme de Venn}

\begin{exo}[Compléter un diagramme de Venn]
\cas{1)}
\begin{minipage}[c]{0.58\linewidth}\raggedright
On part de la zone commune : $12$ bénévoles font les deux.
\begin{calculs}
35-12 &= \res{23} &&\justif{réparation seulement}\\
28-12 &= \res{16} &&\justif{aide aux devoirs seulement}\\
60-(23+12+16) &= \res{9} &&\justif{ni l'un ni l'autre : l'accueil}
\end{calculs}
\end{minipage}\hfill
\begin{minipage}[c]{0.38\linewidth}\centering
\begin{tikzpicture}[scale=0.8]
  \venn{R}{D}{$23$}{$12$}{$16$}{$9$}
\end{tikzpicture}
\end{minipage}
\cas{2)}
\begin{calculs}
23+12+16 &= \res{51} &&\justif{au moins une activité : les trois zones des cercles}\\
23+16 &= \res{39} &&\justif{une seule : les deux zones \og seulement\fg{}}
\end{calculs}
\explique{$51$ bénévoles font au moins une des deux activités, \res{$9$} n'en font aucune
(ils tiennent l'accueil) et $39$ en font une seule.}{aucune : la zone hors des cercles}
\rem{Erreur fréquente : additionner $35+28=63$, plus que $60$ ; les $12$ bénévoles qui font
les deux sont alors comptés deux fois.}
\end{exo}

\begin{exo}[Partir de la zone commune]
\cas{1)}
\begin{minipage}[c]{0.58\linewidth}\raggedright
On part de la zone commune : $8$ élèves font les deux.
\begin{calculs}
18-8 &= \res{10} &&\justif{sport seulement}\\
30-5 &= 25 &&\justif{élèves dans au moins un des cercles}\\
25-(10+8) &= \res{7} &&\justif{musique seulement}
\end{calculs}
\end{minipage}\hfill
\begin{minipage}[c]{0.38\linewidth}\centering
\begin{tikzpicture}[scale=0.8]
  \venn{S}{M}{$10$}{$8$}{$7$}{$5$}
\end{tikzpicture}
\end{minipage}
\cas{2)}
\begin{calculs}
8+7 &= \res{15} &&\justif{les deux zones du cercle $M$}
\end{calculs}
\explique{$15$ élèves jouent d'un instrument de musique.}{zone commune et \og musique
seulement\fg{}}
\cas{3)}
\begin{calculs}
10+7 &= \res{17} &&\justif{les deux zones \og seulement\fg{}}
\end{calculs}
\explique{$17$ élèves font une seule des deux activités.}{la zone commune n'est pas
comptée}
\end{exo}

\begin{exo}[Ranger des nombres]
\cas{a)}
\begin{minipage}[c]{0.5\linewidth}\raggedright
$A=\ens{2\,;4\,;6\,;8\,;10\,;12}$ et $B=\ens{3\,;6\,;9\,;12}$.

\smallskip
Zone commune : $A\cap B=\res{\ens{6\,;12}}$ ; les autres nombres de $A$ vont dans \og $A$
seulement\fg{}, ceux de $B$ dans \og $B$ seulement\fg{}, et les nombres restants hors des
cercles.
\end{minipage}\hfill
\begin{minipage}[c]{0.46\linewidth}\centering
\begin{tikzpicture}[scale=1.3]
  \vennliste{A}{B}{$2$\enspace$4$\\$8$\enspace$10$}{$6$\\$12$}{$3$\\$9$}%
    {$1$ \; $5$ \; $7$ \; $11$}
\end{tikzpicture}
\end{minipage}
\cas{b)}
\begin{minipage}[c]{0.5\linewidth}\raggedright
$C=\ens{4\,;8\,;12}$ : tous les multiples de $4$ sont pairs, donc tous les nombres de $C$
sont dans la zone commune : $A\cap C=\res{\ens{4\,;8\,;12}}$.
\end{minipage}\hfill
\begin{minipage}[c]{0.46\linewidth}\centering
\begin{tikzpicture}[scale=1.3]
  \vennliste{A}{C}{$2$\\$6$\\$10$}{$4$\\$8$\\$12$}{}%
    {$1$ \; $3$ \; $5$ \; $7$ \; $9$ \; $11$}
\end{tikzpicture}
\end{minipage}
\cas{c)}
\begin{minipage}[c]{0.5\linewidth}\raggedright
$D=\ens{1\,;3\,;5\,;7\,;9\,;11}$ : aucun nombre n'est à la fois pair et impair, donc
$A\cap D=\res{\varnothing}$ ; et tout nombre est pair ou impair, donc aucun nombre n'est
hors des cercles.
\end{minipage}\hfill
\begin{minipage}[c]{0.46\linewidth}\centering
\begin{tikzpicture}[scale=1.3]
  \vennliste{A}{D}{$2$\enspace$4$\\$6$\enspace$8$\\$10$\enspace$12$}{}%
    {$1$\enspace$3$\\$5$\enspace$7$\\$9$\enspace$11$}{aucun nombre}
\end{tikzpicture}
\end{minipage}
\cas{Dernières questions}
\explique{Dans le cas c), la zone commune est vide : \res{$A$ et $D$ sont
incompatibles}.}{$A\cap D=\varnothing$ ; on peut aussi dessiner deux cercles séparés}
\explique{Dans le cas b), \res{la zone \og $C$ seulement\fg{} est vide} : tout multiple de
$4$ est pair, donc tout nombre de $C$ est aussi dans $A$.}{$4=2\times 2$}
\end{exo}

\begin{exo}[Lire un diagramme de Venn]
\cas{1)}
\begin{calculs}
70+30 &= \res{100} &&\justif{les deux zones du cercle $V$}
\end{calculs}
\explique{$100$ visiteurs ont essayé un vélo électrique.}{zone commune comprise}
\cas{2)}
\explique{$30$ est le \res{nombre de visiteurs qui ont essayé le vélo électrique
\textbf{et} la trottinette électrique}.}{la zone commune, $V\cap T$}
\explique{$60$ est le \res{nombre de visiteurs qui n'ont essayé ni le vélo ni la
trottinette}.}{hors des deux cercles}
\cas{3)}
\begin{calculs}
70+40 &= \res{110} &&\justif{les deux zones \og seulement\fg{}}
\end{calculs}
\explique{$110$ visiteurs ont essayé un seul des deux engins.}{la zone commune n'est pas
comptée}
\cas{4)}
\explique{$V\cup T$ est l'ensemble des \res{visiteurs qui ont essayé le vélo électrique ou
la trottinette électrique} (au moins un des deux).}{la réunion}
\begin{calculs}
70+30+40 &= \res{140} &&\justif{les trois zones des cercles ; ou $200-60$}
\end{calculs}
\end{exo}

\partie{Équiprobabilité}

\begin{exo}[Calculer des probabilités]
\cas{a)}
\explique{Les boules sont indiscernables : il y a équiprobabilité, avec $30$ issues.}{chaque
boule a la même chance}
\begin{calculs}
P(A) &= \dfrac{10}{30} &&\justif{$A=\ens{3\,;6\,;9\,;\dots\,;30}$ : $10$ multiples de $3$}\\
P(A) &= \dfrac{\barre{10}\times 1}{\barre{10}\times 3} &&\justif{on simplifie par $10$}\\
P(A) &= \res{\dfrac{1}{3}} &&\\
P(B) &= \dfrac{6}{30} = \res{\dfrac{1}{5}} &&\justif{$B=\ens{5\,;10\,;15\,;20\,;25\,;30}$ ;
on simplifie par $6$}\\
P(A\cap B) &= \dfrac{2}{30} = \res{\dfrac{1}{15}} &&\justif{$A\cap B=\ens{15\,;30}$ ; on
simplifie par $2$}
\end{calculs}
\cas{b)}
\explique{Le bénévole est choisi au hasard : il y a équiprobabilité, et on divise chaque
effectif par $60$.}{$60$ bénévoles en tout}
\begin{calculs}
P(R\cap D) &= \dfrac{12}{60} = \res{0{,}2} &&\justif{la zone commune}\\
P(R\cup D) &= \dfrac{23+12+16}{60} &&\justif{les trois zones des cercles}\\
P(R\cup D) &= \dfrac{51}{60} = \res{0{,}85} &&\\
P(\overline{R}\cap\overline{D}) &= \dfrac{9}{60} = \res{0{,}15} &&\justif{ni réparation ni
aide aux devoirs}
\end{calculs}
\end{exo}

\begin{exo}[Dans la boîte à outils]
\explique{La vis est choisie au hasard : il y a équiprobabilité, et on divise chaque
effectif par $100$.}{$35+15+5+10+20+15=100$ vis}
\cas{1)}
\begin{calculs}
P(C) &= \dfrac{10+20+15}{100} &&\justif{$C$ : \og tête de $5$ mm\fg{} ; toute la deuxième
ligne}\\
P(C) &= \dfrac{45}{100} = \res{0{,}45} &&
\end{calculs}
\cas{2)}
\begin{calculs}
P(L) &= \dfrac{15+5+20+15}{100} &&\justif{$L$ : \og plus de $30$ mm\fg{} ; colonnes $40$ et
$50$ mm}\\
P(L) &= \dfrac{55}{100} = \res{0{,}55} &&
\end{calculs}
\cas{3)}
\begin{calculs}
P(E) &= \dfrac{35}{100} = \res{0{,}35} &&\justif{$E$ : \og tête de $4$ mm et $30$ mm de
long\fg{} : une seule case}
\end{calculs}
\cas{4)}
\begin{calculs}
P(F) &= \dfrac{5+15}{100} &&\justif{$F$ : \og $50$ mm de long\fg{} ; toute la dernière
colonne}\\
P(F) &= \dfrac{20}{100} = \res{0{,}2} &&
\end{calculs}
\rem{On nomme chaque événement par une lettre : on écrit ensuite $P(C)$ plutôt qu'une
phrase.}
\end{exo}

\begin{exo}[Vrai ou faux]
\cas{a)}
\explique{\res{Vrai} : obtenir $7$ est impossible avec un dé numéroté de $1$ à $6$, et un
événement impossible a une probabilité égale à $0$.}{définition 5}
\cas{b)}
\explique{\res{Faux} : une probabilité est toujours comprise entre $0$ et $1$.}{$1{,}2$ est
plus grand que $1$}
\cas{c)}
\explique{\res{Faux} : l'urne contient $5+3=8$ boules, dont $3$ vertes ; la probabilité de
tirer une boule verte est $\res{\dfrac{3}{8}}$.}{on divise par le nombre \textbf{total} de
boules, pas par le nombre de boules rouges}
\cas{d)}
\explique{\res{Faux} : le $2$ est écrit sur trois faces sur six, il a une probabilité de
$\dfrac{3}{6}$, soit $\dfrac{1}{2}$ ; le $6$, sur une seule face, a une probabilité de
$\dfrac{1}{6}$.}{les six faces sont équiprobables, pas les trois nombres lus}
\end{exo}

\begin{exo}[Un jeu de 32 cartes]
\explique{La carte est tirée au hasard : il y a équiprobabilité, avec $32$ issues.}{$4$
couleurs de $8$ cartes}
\cas{1)}
\begin{calculs}
P(R) &= \dfrac{4}{32} = \res{\dfrac{1}{8}} &&\justif{un roi par couleur ; on simplifie par
$4$}\\
P(C) &= \dfrac{8}{32} = \res{\dfrac{1}{4}} &&\justif{$8$ cartes de cœur ; on simplifie par
$8$}\\
P(F) &= \dfrac{12}{32} = \res{\dfrac{3}{8}} &&\justif{$3$ figures par couleur ; on simplifie
par $4$}
\end{calculs}
\cas{2)}
\explique{$R\cap C$ : \og \res{tirer le roi de cœur}\fg{}.}{roi \textbf{et} cœur}
\begin{calculs}
P(R\cap C) &= \res{\dfrac{1}{32}} &&\justif{une seule carte}
\end{calculs}
\cas{3)}
\explique{Une figure de cœur est le valet, la dame ou le roi de cœur : $3$ cartes.}{c'est
l'événement $F\cap C$}
\begin{calculs}
P(F\cap C) &= \res{\dfrac{3}{32}} &&
\end{calculs}
\end{exo}

\partie{Loi de probabilité}

\begin{exo}[Utiliser une loi de probabilité]
\cas{1)}
\begin{calculs}
0{,}05+0{,}10+0{,}15+p+0{,}30+0{,}15 &= 1 &&\justif{la somme des probabilités vaut $1$}\\
0{,}75+p &= 1 &&\\
p &= 1-0{,}75 &&\\
p &= \res{0{,}25} &&
\end{calculs}
\cas{2)}
\begin{calculs}
P(A) &= 0{,}25+0{,}30+0{,}15 &&\justif{$A=\ens{6\,;8\,;10}$ : on additionne}\\
P(A) &= \res{0{,}70} &&
\end{calculs}
\cas{3)}
\begin{calculs}
P(B) &= 0{,}05+0{,}10+0{,}15 &&\justif{$B=\ens{0\,;2\,;4}$}\\
P(B) &= \res{0{,}30} &&
\end{calculs}
\rem{On peut vérifier : $A$ et $B$ réunissent toutes les issues, sans issue commune, et
$0{,}70+0{,}30=1$.}
\end{exo}

\begin{exo}[À la borne de recharge]
\cas{1)}
\begin{calculs}
0{,}35+0{,}40+0{,}20+p &= 1 &&\justif{la somme des probabilités vaut $1$}\\
0{,}95+p &= 1 &&\\
p &= 1-0{,}95 &&\\
p &= \res{0{,}05} &&
\end{calculs}
\cas{2)}
\begin{calculs}
0{,}35+0{,}40 &= \res{0{,}75} &&\justif{\og moins de $5$ min\fg{} et \og de $5$ à $10$
min\fg{}}
\end{calculs}
\explique{La probabilité d'attendre moins de $10$ minutes est $0{,}75$.}{on additionne les
deux premières lignes}
\cas{3)}
\begin{calculs}
0{,}40+0{,}20+0{,}05 &= \res{0{,}65} &&\justif{toutes les lignes sauf \og moins de $5$
min\fg{}}
\end{calculs}
\cas{4)}
\explique{Le contraire de \og attendre moins de $10$ minutes\fg{} est \og \res{attendre $10$
minutes ou plus}\fg{}.}{les deux dernières lignes du tableau}
\begin{calculs}
0{,}20+0{,}05 &= \res{0{,}25} &&\justif{on retrouve $1-0{,}75$}
\end{calculs}
\end{exo}

\begin{exo}[Un dé truqué]
\cas{1)}
\begin{calculs}
0{,}05+4p+0{,}35 &= 1 &&\justif{les quatre autres faces ont chacune la probabilité $p$}\\
0{,}40+4p &= 1 &&\\
4p &= 1-0{,}40 &&\\
4p &= 0{,}60 &&\\
p &= \dfrac{0{,}60}{4} &&\justif{on divise par $4$}\\
p &= \res{0{,}15} &&
\end{calculs}
\cas{2)}
\begin{calculs}
P(A) &= 0{,}05+0{,}15+0{,}15 &&\justif{$A=\ens{1\,;3\,;5}$}\\
P(A) &= \res{0{,}35} &&
\end{calculs}
\cas{3)}
\begin{calculs}
P(B) &= 0{,}15+0{,}15+0{,}15+0{,}35 &&\justif{$B=\ens{3\,;4\,;5\,;6}$}\\
P(B) &= \res{0{,}80} &&
\end{calculs}
\cas{4)}
\explique{Avec un dé équilibré, il y aurait équiprobabilité : $P(A)=\dfrac{3}{6}$, soit
$\res{\dfrac{1}{2}}$.}{$3$ faces impaires sur $6$}
\rem{Le dé truqué donne un nombre impair moins souvent : $0{,}35$ au lieu de $0{,}5$.}
\end{exo}

\partie{Événement contraire, réunion}

\begin{exo}[Contraire et réunion]
\cas{1)}
\begin{calculs}
P(\overline{A}) &= 1-P(A) &&\justif{probabilité du contraire}\\
P(\overline{A}) &= 1-0{,}25 = \res{0{,}75} &&
\end{calculs}
\explique{$75\,\%$ des colis ne sont pas volumineux.}{$\overline{A}$ : \og le colis n'est pas
volumineux\fg{}}
\cas{2)}
\begin{calculs}
P(A\cup B) &= P(A)+P(B)-P(A\cap B) &&\justif{\og volumineux ou fragile\fg{} : la réunion}\\
P(A\cup B) &= 0{,}25+0{,}15-0{,}05 &&\\
P(A\cup B) &= \res{0{,}35} &&
\end{calculs}
\cas{3)}
\begin{calculs}
P(\overline{A}\cap\overline{B}) &= 1-P(A\cup B) &&\justif{\og ni l'un ni l'autre\fg{} :
contraire de \og au moins un\fg{}}\\
P(\overline{A}\cap\overline{B}) &= 1-0{,}35 = \res{0{,}65} &&
\end{calculs}
\cas{4)}
\begin{calculs}
P(A\cup C) &= P(A)+P(C) &&\justif{$A$ et $C$ sont incompatibles : une enveloppe n'est
jamais volumineuse}\\
P(A\cup C) &= 0{,}25+0{,}3 = \res{0{,}55} &&
\end{calculs}
\end{exo}

\begin{exo}[Retrouver une intersection]
\cas{1)}
\begin{calculs}
P(\overline{B}) &= 1-P(B) &&\justif{probabilité du contraire}\\
P(\overline{B}) &= 1-0{,}35 = \res{0{,}65} &&
\end{calculs}
\cas{2)}
\begin{calculs}
P(B\cup C) &= P(B)+P(C) &&\justif{$B$ et $C$ sont incompatibles}\\
P(B\cup C) &= 0{,}35+0{,}4 = \res{0{,}75} &&
\end{calculs}
\cas{3)}
\explique{On part de la formule de la réunion : la seule valeur inconnue est
$P(A\cap B)$.}{on connaît $P(A)$, $P(B)$ et $P(A\cup B)$}
\begin{calculs}
P(A\cup B) &= P(A)+P(B)-P(A\cap B) &&\justif{formule de la réunion}\\
0{,}45 &= 0{,}2+0{,}35-P(A\cap B) &&\justif{on remplace par les valeurs connues}\\
0{,}45 &= 0{,}55-P(A\cap B) &&\\
P(A\cap B) &= 0{,}55-0{,}45 &&\justif{on isole $P(A\cap B)$}\\
P(A\cap B) &= \res{0{,}1} &&
\end{calculs}
\cas{4)}
\begin{calculs}
P(\overline{A}\cap\overline{B}) &= 1-P(A\cup B) &&\justif{contraire de \og au moins un des
deux\fg{}}\\
P(\overline{A}\cap\overline{B}) &= 1-0{,}45 = \res{0{,}55} &&
\end{calculs}
\end{exo}

\begin{exo}[Des données impossibles]
\cas{a)}
\explique{$P(\overline{A})$ devrait être égale à $1-0{,}6$, soit $0{,}4$, et non $0{,}5$ :
\res{impossible}.}{$P(A)+P(\overline{A})$ vaut toujours $1$}
\cas{b)}
\explique{$A$ et $B$ sont incompatibles, donc $P(A\cup B)=0{,}7+0{,}4=1{,}1$, plus grand que
$1$ : \res{impossible}.}{une probabilité ne dépasse jamais $1$}
\cas{c)}
\explique{Les issues de $A\cap B$ sont toutes dans $A$, donc $P(A\cap B)$ ne peut pas
dépasser $P(A)=0{,}3$ ; or $0{,}5$ est plus grand que $0{,}3$ : \res{impossible}.}{$A\cap B$
fait partie de $A$}
\cas{d)}
\explique{Les issues de $B$ sont toutes dans $A\cup B$, donc $P(A\cup B)$ est au moins égale
à $P(B)=0{,}6$ ; or $0{,}4$ est plus petit que $0{,}6$ : \res{impossible}.}{la réunion
contient chacun des deux événements}
\rem{Autre méthode pour d) : la formule de la réunion donnerait $P(A\cap B)=0{,}5+0{,}6-0{,}4$,
soit $0{,}7$, plus grand que $P(A)=0{,}5$.}
\end{exo}

\begin{exo}[Des roulements à billes]
\cas{1)}
\begin{minipage}[c]{0.58\linewidth}\raggedright
\begin{calculs}
0{,}03\times 5\,000 &= \res{150} &&\justif{roulements de $A$ : $3\,\%$ de $5\,000$}\\
0{,}05\times 5\,000 &= \res{250} &&\justif{roulements de $B$}\\
0{,}01\times 5\,000 &= \res{50} &&\justif{roulements de $A\cap B$}
\end{calculs}
On part de la zone commune : $50$ roulements.
\begin{calculs}
150-50 &= 100 &&\justif{diamètre seulement}\\
250-50 &= 200 &&\justif{surface seulement}\\
5\,000-(100+50+200) &= 4\,650 &&\justif{aucun défaut}
\end{calculs}
\end{minipage}\hfill
\begin{minipage}[c]{0.38\linewidth}\centering
\begin{tikzpicture}[scale=0.8]
  \venn{A}{B}{$100$}{$50$}{$200$}{\footnotesize$4\,650$}
\end{tikzpicture}
\end{minipage}
\cas{2)}
\explique{Le roulement est choisi au hasard dans le lot : il y a équiprobabilité ; on
cherche la probabilité de \og au moins un des deux défauts\fg{}, la réunion.}{$5\,000$
roulements}
\begin{calculs}
P(A\cup B) &= P(A)+P(B)-P(A\cap B) &&\justif{formule de la réunion}\\
P(A\cup B) &= 0{,}03+0{,}05-0{,}01 &&\\
P(A\cup B) &= \res{0{,}07} &&
\end{calculs}
\rem{Avec le diagramme : $\dfrac{100+50+200}{5\,000}=\dfrac{350}{5\,000}=0{,}07$.}
\cas{3)}
\begin{calculs}
P(\overline{A}\cap\overline{B}) &= 1-P(A\cup B) &&\justif{\og aucun défaut\fg{} : contraire
de \og au moins un\fg{}}\\
P(\overline{A}\cap\overline{B}) &= 1-0{,}07 = \res{0{,}93} &&
\end{calculs}
\cas{4)}
\explique{Un seul défaut : les deux zones \og seulement\fg{}, soit $100+200=300$
roulements.}{la zone commune n'est pas comptée}
\begin{calculs}
\dfrac{300}{5\,000} &= \res{0{,}06} &&\justif{on divise par l'effectif total}
\end{calculs}
\end{exo}

\partie{Arbre des possibles}

\begin{exo}[Dénombrer avec un arbre]
\cas{1)}
\begin{minipage}[c]{0.58\linewidth}\raggedright
Premier chiffre : $4$ branches ; au bout de chacune, deuxième chiffre : $4$ branches (un
chiffre peut revenir).
\begin{calculs}
4\times 4 &= \res{16} &&\justif{$4$ choix, puis $4$ choix}
\end{calculs}
Il y a $16$ codes possibles.
\end{minipage}\hfill
\begin{minipage}[c]{0.38\linewidth}\centering
\begin{tikzpicture}[x=1cm,y=0.34cm,font=\small]
  \node (O) at (0,0) {};
  \foreach \ca/\ya in {1/6,2/2,3/-2,4/-6}{%
    \node[evt] (N\ca) at (1.2,\ya) {$\ca$};
    \draw (O.center) -- (N\ca);
    \foreach \cb/\db in {1/1.5,2/0.5,3/-0.5,4/-1.5}{%
      \node[evt,font=\footnotesize] (N\ca\cb) at (2.4,\ya+\db) {$\cb$};
      \draw (N\ca) -- (N\ca\cb);
      \node[right,font=\footnotesize] at (N\ca\cb.east) {\ca\cb};}}
\end{tikzpicture}
\end{minipage}
\cas{2)}
\explique{Un seul code est le bon, parmi $16$ codes qui ont la même chance : la probabilité
de le trouver est $\res{\dfrac{1}{16}}$.}{code choisi au hasard : équiprobabilité}
\cas{3)}
\begin{calculs}
P(A) &= \dfrac{4}{16} = \res{\dfrac{1}{4}} &&\justif{$A=\ens{11\,;22\,;33\,;44}$ ; on
simplifie par $4$}
\end{calculs}
\cas{4)}
\begin{minipage}[c]{0.58\linewidth}\raggedright
Premier jeton : $3$ branches ; il n'est pas remis dans le sac, donc il reste $2$ jetons pour
le deuxième : $2$ branches au bout de chacune.
\begin{calculs}
3\times 2 &= \res{6} &&\justif{$3$ choix, puis $2$ choix}
\end{calculs}
Il y a $6$ tirages possibles.
\end{minipage}\hfill
\begin{minipage}[c]{0.38\linewidth}\centering
\begin{tikzpicture}[x=1cm,y=0.4cm,font=\small]
  \node (O) at (0,0) {};
  \foreach \ca/\ya/\reste in {R/3/{V,J},V/0/{R,J},J/-3/{R,V}}{%
    \node[evt] (N\ca) at (1.2,\ya) {\ca};
    \draw (O.center) -- (N\ca);
    \foreach \cb [count=\k] in \reste {%
      \node[evt] (N\ca\cb) at (2.4,\ya+1.5-\k) {\cb};
      \draw (N\ca) -- (N\ca\cb);
      \node[right,font=\footnotesize] at (N\ca\cb.east) {\ca\cb};}}
\end{tikzpicture}
\end{minipage}
\cas{5)}
\begin{calculs}
P(B) &= \dfrac{2}{6} = \res{\dfrac{1}{3}} &&\justif{$B=\ens{\text{RV}\,;\text{RJ}}$ ; on
simplifie par $2$}
\end{calculs}
\cas{6)}
\begin{calculs}
P(C) &= \dfrac{4}{6} = \res{\dfrac{2}{3}} &&\justif{$C=\ens{\text{RV}\,;\text{RJ}\,;\text{VR}\,;\text{JR}}$
; on simplifie par $2$}
\end{calculs}
\rem{$B$ et $C$ sont différents : dans $C$, le jeton rouge peut sortir en premier ou en
second.}
\end{exo}

\begin{exo}[Trois lancers]
\cas{1)}
\begin{minipage}[c]{0.58\linewidth}\raggedright
Trois lancers : $2$ branches au bout de chaque nœud, à chaque étape.
\begin{calculs}
2\times 2\times 2 &= \res{8} &&\justif{$2$ choix à chacun des trois lancers}
\end{calculs}
Il y a $8$ issues.
\end{minipage}\hfill
\begin{minipage}[c]{0.38\linewidth}\centering
\begin{tikzpicture}[x=1cm,y=0.42cm,font=\small]
  \node (O) at (0,0) {};
  \foreach \ca/\ya in {P/2,F/-2}{%
    \node[evt] (N\ca) at (1.1,\ya) {\ca};
    \draw (O.center) -- (N\ca);
    \foreach \cb/\db in {P/1,F/-1}{%
      \node[evt] (N\ca\cb) at (2.2,\ya+\db) {\cb};
      \draw (N\ca) -- (N\ca\cb);
      \foreach \cc/\dc in {P/0.5,F/-0.5}{%
        \node[evt,font=\footnotesize] (N\ca\cb\cc) at (3.3,\ya+\db+\dc) {\cc};
        \draw (N\ca\cb) -- (N\ca\cb\cc);
        \node[right,font=\footnotesize] at (N\ca\cb\cc.east) {\ca\cb\cc};}}}
\end{tikzpicture}
\end{minipage}
\cas{2)}
\explique{Un seul chemin donne PPP, parmi $8$ issues qui ont la même chance : la probabilité
est $\res{\dfrac{1}{8}}$.}{pièce équilibrée : équiprobabilité}
\cas{3)}
\begin{calculs}
P(A) &= \res{\dfrac{3}{8}} &&\justif{$A=\ens{\text{PFF}\,;\text{FPF}\,;\text{FFP}}$ : $3$
issues sur $8$}
\end{calculs}
\cas{4)}
\explique{On note $E$ \og obtenir au moins une fois Pile\fg{} ; son contraire est \og ne
jamais obtenir Pile\fg{} : $\overline{E}=\ens{\text{FFF}}$.}{une seule issue, plus simple à
compter}
\begin{calculs}
P(E) &= 1-P(\overline{E}) &&\justif{probabilité du contraire}\\
P(E) &= 1-\dfrac{1}{8} = \res{\dfrac{7}{8}} &&
\end{calculs}
\end{exo}

\begin{exo}[La tenue du cycliste]
\cas{1)}
\begin{minipage}[c]{0.58\linewidth}\raggedright
Maillot : $3$ branches ; au bout de chacune, casque : $2$ branches.
\begin{calculs}
3\times 2 &= \res{6} &&\justif{$3$ choix, puis $2$ choix}
\end{calculs}
Il y a $6$ tenues possibles.
\end{minipage}\hfill
\begin{minipage}[c]{0.38\linewidth}\centering
\begin{tikzpicture}[x=1cm,y=0.42cm,font=\footnotesize]
  \node (O) at (0,0) {};
  \foreach \ca/\ya in {rouge/3,bleu/0,blanc/-3}{%
    \node[evt] (N\ca) at (1.4,\ya) {\ca};
    \draw (O.center) -- (N\ca);
    \foreach \cb/\db in {noir/0.7,blanc/-0.7}{%
      \node[evt] (N\ca\cb) at (3,\ya+\db) {\cb};
      \draw (N\ca) -- (N\ca\cb);}}
  \node[font=\scriptsize] at (1.4,4.4) {maillot};
  \node[font=\scriptsize] at (3,4.4) {casque};
\end{tikzpicture}
\end{minipage}
\cas{2)}
\explique{Une seule tenue est toute blanche, parmi $6$ tenues qui ont la même chance : la
probabilité est $\res{\dfrac{1}{6}}$.}{choix au hasard : équiprobabilité}
\cas{3)}
\begin{calculs}
P(M) &= \dfrac{2}{6} = \res{\dfrac{1}{3}} &&\justif{maillot blanc : $2$ tenues (casque noir
ou blanc)}\\
P(C) &= \dfrac{3}{6} = \res{\dfrac{1}{2}} &&\justif{casque blanc : $3$ tenues, une par
maillot}\\
P(M\cap C) &= \dfrac{1}{6} &&\justif{question 2}
\end{calculs}
\begin{calculs}
P(M\cup C) &= P(M)+P(C)-P(M\cap C) &&\justif{formule de la réunion}\\
P(M\cup C) &= \dfrac{2}{6}+\dfrac{3}{6}-\dfrac{1}{6} &&\justif{on garde le dénominateur
$6$}\\
P(M\cup C) &= \dfrac{4}{6} = \res{\dfrac{2}{3}} &&\justif{on simplifie par $2$}
\end{calculs}
\rem{On vérifie sur l'arbre : $4$ tenues sur $6$ ont au moins un élément blanc.}
\end{exo}

\partie{Probabilité conditionnelle et tableau croisé}

\begin{exo}[Probabilités dans un tableau croisé]
\cas{1)}
\begin{minipage}[c]{0.5\linewidth}\raggedright
\begin{calculs}
100+100 &= \res{200} &&\justif{total du matin}\\
400-200 &= \res{200} &&\justif{total du soir}\\
200-50 &= \res{150} &&\justif{annuelles le soir}\\
100+50 &= \res{150} &&\justif{total des mensuelles}\\
100+150 &= \res{250} &&\justif{total des annuelles}
\end{calculs}
\end{minipage}\hfill
\begin{minipage}[c]{0.46\linewidth}\centering
\renewcommand{\arraystretch}{1.3}
\begin{tabular}{|l|c|c|c|}
\hline
 & Mensuelle & Annuelle & Total\\
\hline
Matin & $100$ & $100$ & $200$\\
\hline
Soir & $50$ & $150$ & $200$\\
\hline
Total & $150$ & $250$ & $400$\\
\hline
\end{tabular}
\end{minipage}
\cas{2)}
\begin{calculs}
P(A) &= \dfrac{250}{400} = \res{0{,}625} &&\justif{$250$ formules annuelles sur $400$
abonnés}\\
P(A\cap S) &= \dfrac{150}{400} = \res{0{,}375} &&\justif{annuelles \textbf{et} le soir,
sur le total}
\end{calculs}
\cas{3)}
\begin{calculs}
P_A(S) &= \dfrac{150}{250} = \res{0{,}6} &&\justif{on ne garde que la colonne
\og Annuelle\fg{}}
\end{calculs}
\explique{$60\,\%$ des abonnés à l'année viennent le soir.}{\og parmi les formules
annuelles\fg{}}
\cas{4)}
\begin{calculs}
P_S(A) &= \dfrac{150}{200} = \res{0{,}75} &&\justif{on ne garde que la ligne \og Soir\fg{}}
\end{calculs}
\explique{$75\,\%$ des abonnés qui viennent le soir ont une formule annuelle.}{$P_A(S)$ et
$P_S(A)$ ne sont pas égales}
\end{exo}

\begin{exo}[Une équipe de rugby]
\cas{1)}
\begin{minipage}[c]{0.46\linewidth}\raggedright
\begin{calculs}
24-18 &= \res{6} &&\justif{avants de $1{,}85$ m ou moins}\\
16-2 &= \res{14} &&\justif{arrières de $1{,}85$ m ou moins}\\
18+2 &= \res{20} &&\justif{joueurs de plus de $1{,}85$ m}
\end{calculs}
\end{minipage}\hfill
\begin{minipage}[c]{0.5\linewidth}\centering
\renewcommand{\arraystretch}{1.3}
\begin{tabular}{|l|c|c|c|}
\hline
 & $G$ & $\overline{G}$ & Total\\
\hline
Avants ($A$) & $18$ & $6$ & $24$\\
\hline
Arrières & $2$ & $14$ & $16$\\
\hline
Total & $20$ & $20$ & $40$\\
\hline
\end{tabular}
\end{minipage}
\cas{2)}
\explique{Le joueur est choisi au hasard : il y a équiprobabilité, et on divise par
$40$.}{$40$ joueurs}
\begin{calculs}
P(A) &= \dfrac{24}{40} = \res{0{,}6} &&\justif{la ligne des avants}\\
P(G) &= \dfrac{20}{40} = \res{0{,}5} &&\justif{la colonne $G$}\\
P(A\cap G) &= \dfrac{18}{40} = \res{0{,}45} &&\justif{avants \textbf{et} plus de $1{,}85$ m}
\end{calculs}
\cas{3)}
\begin{calculs}
P_A(G) &= \dfrac{18}{24} = \res{0{,}75} &&\justif{on ne garde que la ligne \og Avants\fg{} :
$24$ joueurs}
\end{calculs}
\explique{$75\,\%$ des avants mesurent plus de $1{,}85$ m.}{\og parmi les avants\fg{}}
\cas{4)}
\begin{calculs}
P_G(A) &= \dfrac{18}{20} = \res{0{,}9} &&\justif{on ne garde que la colonne $G$ : $20$
joueurs}
\end{calculs}
\explique{$90\,\%$ des joueurs de plus de $1{,}85$ m sont des avants.}{\og parmi les joueurs
de plus de $1{,}85$ m\fg{}}
\cas{5)}
\explique{$P(A\cap G)=0{,}45$ : \res{parmi \textbf{tous} les joueurs, $45\,\%$ sont des
avants de plus de $1{,}85$ m}. $P_A(G)=0{,}75$ : \res{parmi les \textbf{avants} seulement,
$75\,\%$ mesurent plus de $1{,}85$ m}.}{même numérateur, $18$ ; dénominateur $40$ ou $24$}
\end{exo}

\begin{exo}[Au magasin de cycles]
\cas{1)}
\begin{minipage}[c]{0.5\linewidth}\raggedright
\begin{calculs}
0{,}55\times 400 &= 220 &&\justif{VTT}\\
0{,}30\times 400 &= 120 &&\justif{VTC}\\
400-(220+120) &= 60 &&\justif{vélos de ville}\\
0{,}60\times 220 &= \res{132} &&\justif{VTT électriques}\\
0{,}50\times 120 &= \res{60} &&\justif{VTC électriques}\\
0{,}80\times 60 &= \res{48} &&\justif{vélos de ville électriques}\\
132+60+48 &= \res{240} &&\justif{total des électriques}
\end{calculs}
Les vélos non électriques s'obtiennent par différence :
$220-132=88$ ; $120-60=60$ ; $60-48=12$ ; $400-240=160$.
\end{minipage}\hfill
\begin{minipage}[c]{0.46\linewidth}\centering
\renewcommand{\arraystretch}{1.3}
\begin{tabular}{|l|c|c|c|}
\hline
 & Électrique & Non électrique & Total\\
\hline
VTT & $132$ & $88$ & $220$\\
\hline
VTC & $60$ & $60$ & $120$\\
\hline
Ville & $48$ & $12$ & $60$\\
\hline
Total & $240$ & $160$ & $400$\\
\hline
\end{tabular}
\end{minipage}
\cas{2)}
\begin{calculs}
P(E) &= \dfrac{240}{400} = \res{0{,}6} &&\justif{$240$ vélos électriques sur $400$}\\
P(\overline{E}) &= 1-0{,}6 = \res{0{,}4} &&\justif{probabilité du contraire}
\end{calculs}
\cas{3)}
\begin{calculs}
P_E(T) &= \dfrac{132}{240} = \res{0{,}55} &&\justif{on ne garde que la colonne
\og Électrique\fg{} : $240$ vélos}
\end{calculs}
\explique{$55\,\%$ des vélos électriques vendus sont des VTT.}{\og parmi les vélos
électriques\fg{}}
\cas{4)}
\explique{On sait que le vélo n'est pas électrique : on ne garde que la colonne \og Non
électrique\fg{}, $160$ vélos, dont $60$ VTC. On note $V$ \og le vélo est un VTC\fg{}.}{probabilité
sachant $\overline{E}$}
\begin{calculs}
P_{\overline{E}}(V) &= \dfrac{60}{160} = \res{0{,}375} &&
\end{calculs}
\end{exo}

\begin{exo}[Vacances sportives]
\cas{1)}
\begin{minipage}[c]{0.58\linewidth}\raggedright
On part de la zone commune : $12$ personnes font les deux.
\begin{calculs}
48-12 &= \res{36} &&\justif{VTT seulement}\\
30-12 &= \res{18} &&\justif{escalade seulement}\\
120-(36+12+18) &= \res{54} &&\justif{ni l'un ni l'autre}
\end{calculs}
\end{minipage}\hfill
\begin{minipage}[c]{0.38\linewidth}\centering
\begin{tikzpicture}[scale=0.8]
  \venn{T}{E}{$36$}{$12$}{$18$}{$54$}
\end{tikzpicture}
\end{minipage}
\cas{2)}
\begin{minipage}[c]{0.5\linewidth}\raggedright
Chaque zone du diagramme devient une case du tableau ; les totaux se complètent :
\begin{calculs}
120-48 &= \res{72} &&\justif{total de $\overline{T}$}\\
120-30 &= \res{90} &&\justif{total de $\overline{E}$}
\end{calculs}
\end{minipage}\hfill
\begin{minipage}[c]{0.46\linewidth}\centering
\renewcommand{\arraystretch}{1.3}
\begin{tabular}{|l|c|c|c|}
\hline
 & $T$ & $\overline{T}$ & Total\\
\hline
$E$ & $12$ & $18$ & $30$\\
\hline
$\overline{E}$ & $36$ & $54$ & $90$\\
\hline
Total & $48$ & $72$ & $120$\\
\hline
\end{tabular}
\end{minipage}
\cas{3)}
\begin{calculs}
P(\overline{T}\cap\overline{E}) &= \dfrac{54}{120} = \res{0{,}45} &&\justif{ni VTT ni
escalade}
\end{calculs}
\cas{4)}
\begin{calculs}
P_T(E) &= \dfrac{12}{48} = \res{0{,}25} &&\justif{on ne garde que la colonne $T$ : $48$
personnes}\\
P_E(T) &= \dfrac{12}{30} = \res{0{,}4} &&\justif{on ne garde que la ligne $E$ : $30$
personnes}
\end{calculs}
\explique{$25\,\%$ des personnes qui font du VTT font aussi de l'escalade ; $40\,\%$ des
personnes qui font de l'escalade font aussi du VTT.}{même numérateur, deux populations
différentes}
\end{exo}

\partie{Formule des probabilités conditionnelles}

\begin{exo}[Traduire et utiliser la formule]
\cas{a)}
\explique{Données : $P(C)=\res{0{,}4}$ et $P(C\cap S)=\res{0{,}1}$.}{\og une pompe à chaleur
\textbf{et} des panneaux solaires\fg{} : une intersection}
\begin{calculs}
P_C(S) &= \dfrac{P(C\cap S)}{P(C)} &&\justif{formule de la propriété 4}\\
P_C(S) &= \dfrac{0{,}1}{0{,}4} &&\\
P_C(S) &= \res{0{,}25} &&
\end{calculs}
\explique{$25\,\%$ des maisons qui ont une pompe à chaleur ont aussi des panneaux
solaires.}{$P_C(S)$ : parmi les maisons avec pompe à chaleur}
\cas{b)}
\explique{Données : $P(D)=\res{0{,}6}$ et $P_D(V)=\res{0{,}3}$.}{\og parmi eux\fg{} :
probabilité sachant $D$}
\begin{calculs}
P(D\cap V) &= P(D)\times P_D(V) &&\justif{formule de la propriété 4}\\
P(D\cap V) &= 0{,}6\times 0{,}3 &&\\
P(D\cap V) &= \res{0{,}18} &&
\end{calculs}
\explique{$18\,\%$ des élèves déjeunent à la cantine et choisissent le menu
végétarien.}{$D$ \textbf{et} $V$, sur l'ensemble des élèves}
\end{exo}

\begin{exo}[Conditionnelle ou pas ?]
\cas{1)}
\explique{a) Pas conditionnelle : $P(R)=\res{0{,}4}$.}{\og dans l'année\fg{} : sur tous les
jours}
\explique{b) Conditionnelle : $P_R(V)=\res{0{,}75}$.}{\og des jours pluvieux\fg{} : parmi les
jours de pluie}
\explique{c) Conditionnelle : $P_{\overline{V}}(R)=\res{0{,}2}$.}{\og parmi les jours sans
vent fort\fg{} : sachant $\overline{V}$}
\explique{d) Pas conditionnelle : $P(R\cap V)=\res{0{,}3}$.}{\og à la fois\fg{} : une
intersection}
\explique{e) Conditionnelle : $P_V(R)=\res{0{,}6}$.}{\og quand le vent souffle fort\fg{} :
sachant $V$ ; $6$ sur $10$}
\cas{2)}
\begin{calculs}
P_R(V) &= \dfrac{P(R\cap V)}{P(R)} &&\justif{formule de la propriété 4, avec a) et d)}\\
P_R(V) &= \dfrac{0{,}3}{0{,}4} &&\\
P_R(V) &= \res{0{,}75} &&
\end{calculs}
\explique{On retrouve la valeur de b) : les trois informations sont cohérentes.}{a), b) et
d) disent la même chose de deux façons}
\cas{3)}
\begin{calculs}
P(\overline{R}) &= 1-0{,}4 = \res{0{,}6} &&\justif{probabilité du contraire}
\end{calculs}
\end{exo}

\begin{exo}[Au club nautique]
\cas{1)}
\explique{$P(V)=\res{0{,}6}$ ; $P(K)=\res{0{,}5}$ ; $P(V\cap K)=\res{0{,}15}$.}{\og font les
deux\fg{} : l'intersection}
\cas{2)}
\begin{calculs}
P_V(K) &= \dfrac{P(V\cap K)}{P(V)} &&\justif{formule de la propriété 4}\\
P_V(K) &= \dfrac{0{,}15}{0{,}6} &&\\
P_V(K) &= \res{0{,}25} &&
\end{calculs}
\explique{$25\,\%$ des adhérents qui font de la voile font aussi du kayak.}{parmi les
adhérents qui font de la voile}
\cas{3)}
\explique{On sait que l'adhérent fait du kayak : on cherche $P_K(V)$.}{\og qui fait du
kayak\fg{} : sachant $K$}
\begin{calculs}
P_K(V) &= \dfrac{P(V\cap K)}{P(K)} &&\justif{formule de la propriété 4}\\
P_K(V) &= \dfrac{0{,}15}{0{,}5} &&\\
P_K(V) &= \res{0{,}3} &&
\end{calculs}
\cas{4)}
\begin{calculs}
P(V\cup K) &= P(V)+P(K)-P(V\cap K) &&\justif{\og au moins une des deux\fg{} : la réunion}\\
P(V\cup K) &= 0{,}6+0{,}5-0{,}15 &&\\
P(V\cup K) &= \res{0{,}95} &&
\end{calculs}
\end{exo}

\partie{Arbre pondéré}

\begin{exo}[Construire un arbre pondéré]
\cas{1)}
\begin{minipage}[c]{0.5\linewidth}\raggedright
Données : $P(D)=0{,}4$ ; $P_D(R)=0{,}3$ ; $P_{\overline{D}}(R)=0{,}8$. Les autres branches
s'en déduisent : branches d'un même nœud, somme $1$.
\begin{calculs}
P(\overline{D}) &= 1-0{,}4 = \res{0{,}6} &&\\
P_D(\overline{R}) &= 1-0{,}3 = \res{0{,}7} &&\\
P_{\overline{D}}(\overline{R}) &= 1-0{,}8 = \res{0{,}2} &&
\end{calculs}
\end{minipage}\hfill
\begin{minipage}[c]{0.46\linewidth}\centering
\begin{tikzpicture}
  \arbre{D}{R}{0{,}4}{0{,}6}{0{,}3}{0{,}7}{0{,}8}{0{,}2}
  \feuilles{D}{R}
\end{tikzpicture}
\end{minipage}
\cas{2)}
\begin{calculs}
P(D\cap R) &= P(D)\times P_D(R) &&\justif{chemin $D$ puis $R$ : on multiplie}\\
P(D\cap R) &= 0{,}4\times 0{,}3 &&\\
P(D\cap R) &= \res{0{,}12} &&
\end{calculs}
\explique{$12\,\%$ des grimpeurs sont des débutants qui réussissent la voie
rouge.}{$D$ \textbf{et} $R$, sur l'ensemble des grimpeurs}
\cas{3)}
\begin{calculs}
P(D\cap\overline{R}) &= P(D)\times P_D(\overline{R}) &&\justif{chemin $D$ puis
$\overline{R}$}\\
P(D\cap\overline{R}) &= 0{,}4\times 0{,}7 &&\\
P(D\cap\overline{R}) &= \res{0{,}28} &&
\end{calculs}
\rem{Erreur fréquente : multiplier par $0{,}8$, la probabilité de réussir \textbf{sans} être
débutant ; on suit le chemin qui part de $D$.}
\end{exo}

\begin{exo}[Le jeu de la fête de quartier]
\cas{1)}
\explique{Les $40$ cartes ont la même chance, et $N$ contient les numéros de $1$ à $10$ :
$P(N)=\dfrac{10}{40}=\res{0{,}25}$.}{équiprobabilité (partie III)}
\explique{Si $N$ est réalisé, on fait tourner la première roue : $6$ secteurs sur $10$
portent une étoile, donc $P_N(E)=\dfrac{6}{10}=\res{0{,}6}$.}{probabilité sachant $N$}
\cas{2)}
\begin{minipage}[c]{0.5\linewidth}\raggedright
Sinon, on tourne la seconde roue : $P_{\overline{N}}(E)=\dfrac{3}{10}=0{,}3$. Les autres
branches : somme $1$.
\begin{calculs}
P(\overline{N}) &= 1-0{,}25 = \res{0{,}75} &&\\
P_N(\overline{E}) &= 1-0{,}6 = \res{0{,}4} &&\\
P_{\overline{N}}(\overline{E}) &= 1-0{,}3 = \res{0{,}7} &&
\end{calculs}
\end{minipage}\hfill
\begin{minipage}[c]{0.46\linewidth}\centering
\begin{tikzpicture}
  \arbre{N}{E}{0{,}25}{0{,}75}{0{,}6}{0{,}4}{0{,}3}{0{,}7}
  \feuilles{N}{E}
\end{tikzpicture}
\end{minipage}
\cas{3)}
\begin{calculs}
P(N\cap E) &= P(N)\times P_N(E) &&\justif{chemin $N$ puis $E$}\\
P(N\cap E) &= 0{,}25\times 0{,}6 &&\\
P(N\cap E) &= \res{0{,}15} &&
\end{calculs}
\cas{4)}
\begin{calculs}
P(\overline{N}\cap E) &= P(\overline{N})\times P_{\overline{N}}(E) &&\justif{chemin
$\overline{N}$ puis $E$}\\
P(\overline{N}\cap E) &= 0{,}75\times 0{,}3 &&\\
P(\overline{N}\cap E) &= \res{0{,}225} &&
\end{calculs}
\end{exo}

\begin{exo}[Compléter et lire un arbre]
\cas{1)}
\begin{minipage}[c]{0.5\linewidth}\raggedright
Branches d'un même nœud : somme $1$.
\begin{calculs}
P(\overline{B}) &= 1-0{,}65 = \res{0{,}35} &&\\
P_B(\overline{H}) &= 1-0{,}4 = \res{0{,}6} &&\\
P_{\overline{B}}(\overline{H}) &= 1-0{,}7 = \res{0{,}3} &&
\end{calculs}
\end{minipage}\hfill
\begin{minipage}[c]{0.46\linewidth}\centering
\begin{tikzpicture}
  \arbre{B}{H}{0{,}65}{0{,}35}{0{,}4}{0{,}6}{0{,}7}{0{,}3}
  \feuilles{B}{H}
\end{tikzpicture}
\end{minipage}
\cas{2)}
\explique{$P(B)=0{,}65$ : \res{$65\,\%$ des élèves rentrent en bus}.}{premier niveau : sur
l'ensemble des élèves}
\explique{$P_{\overline{B}}(H)=0{,}7$ : \res{parmi les élèves qui ne rentrent pas en bus,
$70\,\%$ sont chez eux avant 17~h~30}.}{second niveau : probabilité sachant
$\overline{B}$}
\cas{3)}
\begin{calculs}
P(B\cap\overline{H}) &= P(B)\times P_B(\overline{H}) &&\justif{chemin $B$ puis
$\overline{H}$ : après 17~h~30}\\
P(B\cap\overline{H}) &= 0{,}65\times 0{,}6 &&\\
P(B\cap\overline{H}) &= \res{0{,}39} &&
\end{calculs}
\explique{$39\,\%$ des élèves rentrent en bus et sont chez eux après 17~h~30.}{on traduit
par une phrase}
\end{exo}

\partie{Formule des probabilités totales}

\begin{exo}[Probabilités totales]
\cas{1)}
\explique{Deux chemins mènent à $R$ : $D$ puis $R$, et $\overline{D}$ puis $R$.}{$R$ se
réalise avec un débutant ou avec un grimpeur confirmé}
\begin{calculs}
P(R) &= P(D\cap R)+P(\overline{D}\cap R) &&\justif{formule des probabilités totales}\\
P(R) &= 0{,}4\times 0{,}3+0{,}6\times 0{,}8 &&\justif{produit sur chaque chemin}\\
P(R) &= 0{,}12+0{,}48 &&\\
P(R) &= \res{0{,}6} &&
\end{calculs}
\explique{$60\,\%$ des grimpeurs réussissent la voie rouge.}{on traduit par une phrase}
\cas{2)}
\begin{calculs}
P_R(D) &= \dfrac{P(R\cap D)}{P(R)} &&\justif{on sait que $R$ est réalisé : propriété 4}\\
P_R(D) &= \dfrac{0{,}12}{0{,}6} &&\justif{$P(R\cap D)=P(D\cap R)$, calculé en 1)}\\
P_R(D) &= \res{0{,}2} &&
\end{calculs}
\explique{$20\,\%$ des grimpeurs qui réussissent la voie rouge sont débutants.}{à ne pas
confondre avec $P_D(R)=0{,}3$}
\cas{3)}
\begin{minipage}[c]{0.5\linewidth}\raggedright
Données : $P(F)=0{,}6$ ; $P_F(D)=0{,}02$ ; $P_{\overline{F}}(D)=0{,}07$. Les autres
branches s'en déduisent : $P(\overline{F})=0{,}4$ ; $P_F(\overline{D})=0{,}98$ ;
$P_{\overline{F}}(\overline{D})=0{,}93$.
\end{minipage}\hfill
\begin{minipage}[c]{0.46\linewidth}\centering
\begin{tikzpicture}
  \arbre{F}{D}{0{,}6}{0{,}4}{0{,}02}{0{,}98}{0{,}07}{0{,}93}
\end{tikzpicture}
\end{minipage}
\begin{calculs}
P(D) &= P(F)\times P_F(D)+P(\overline{F})\times P_{\overline{F}}(D) &&\justif{deux chemins
mènent à $D$}\\
P(D) &= 0{,}6\times 0{,}02+0{,}4\times 0{,}07 &&\\
P(D) &= 0{,}012+0{,}028 &&\\
P(D) &= \res{0{,}04} &&
\end{calculs}
\cas{4)}
\begin{calculs}
P_D(F) &= \dfrac{P(D\cap F)}{P(D)} &&\justif{on sait que $D$ est réalisé}\\
P_D(F) &= \dfrac{0{,}012}{0{,}04} &&\justif{$P(D\cap F)=0{,}6\times 0{,}02$, calculé en 3)}\\
P_D(F) &= \res{0{,}3} &&
\end{calculs}
\explique{$30\,\%$ des onduleurs tombés en panne viennent du premier fournisseur.}{il
fournit plus d'onduleurs, mais ils tombent moins souvent en panne}
\end{exo}

\begin{exo}[Le sonar du bateau de pêche]
\cas{1)}
\begin{minipage}[c]{0.5\linewidth}\raggedright
Données : $P(B)=0{,}7$ ; $P_B(S)=0{,}8$ ; $P_{\overline{B}}(S)=0{,}1$. Les autres branches
s'en déduisent : branches d'un même nœud, somme $1$.
\begin{calculs}
P(\overline{B}) &= 1-0{,}7 = \res{0{,}3} &&\\
P_B(\overline{S}) &= 1-0{,}8 = \res{0{,}2} &&\\
P_{\overline{B}}(\overline{S}) &= 1-0{,}1 = \res{0{,}9} &&
\end{calculs}
\end{minipage}\hfill
\begin{minipage}[c]{0.46\linewidth}\centering
\begin{tikzpicture}
  \arbre{B}{S}{0{,}7}{0{,}3}{0{,}8}{0{,}2}{0{,}1}{0{,}9}
  \feuilles{B}{S}
\end{tikzpicture}
\end{minipage}
\cas{2)}
\begin{calculs}
P(B\cap S) &= 0{,}7\times 0{,}8 = \res{0{,}56} &&\justif{chemin $B$ puis $S$}\\
P(\overline{B}\cap S) &= 0{,}3\times 0{,}1 = \res{0{,}03} &&\justif{chemin $\overline{B}$
puis $S$}
\end{calculs}
\cas{3)}
\begin{calculs}
P(S) &= P(B\cap S)+P(\overline{B}\cap S) &&\justif{formule des probabilités totales}\\
P(S) &= 0{,}56+0{,}03 &&\justif{question 2}\\
P(S) &= \res{0{,}59} &&
\end{calculs}
\cas{4)}
\explique{Le sonar se trompe quand il signale un banc alors qu'il n'y en a pas : on cherche
$P_S(\overline{B})$.}{on sait que $S$ est réalisé}
\begin{calculs}
P_S(\overline{B}) &= \dfrac{P(S\cap\overline{B})}{P(S)} &&\justif{propriété 4}\\
P_S(\overline{B}) &= \dfrac{0{,}03}{0{,}59} &&\justif{$P(S\cap\overline{B})=P(\overline{B}\cap
S)$}\\
P_S(\overline{B}) &\approx \res{0{,}05} &&\justif{$0{,}0508\dots$ arrondi au centième}
\end{calculs}
\explique{Quand le sonar signale un banc, il se trompe dans environ $5\,\%$ des cas.}{on
traduit par une phrase}
\end{exo}

\begin{exo}[Aller en train, revenir en bus ?]
\cas{1)}
\begin{minipage}[c]{0.5\linewidth}\raggedright
Données : $P(A)=0{,}6$ ; $P_A(R)=0{,}9$ ; $P_{\overline{A}}(R)=0{,}7$. Les autres branches
s'en déduisent : branches d'un même nœud, somme $1$.
\begin{calculs}
P(\overline{A}) &= 1-0{,}6 = \res{0{,}4} &&\\
P_A(\overline{R}) &= 1-0{,}9 = \res{0{,}1} &&\\
P_{\overline{A}}(\overline{R}) &= 1-0{,}7 = \res{0{,}3} &&
\end{calculs}
\end{minipage}\hfill
\begin{minipage}[c]{0.46\linewidth}\centering
\begin{tikzpicture}
  \arbre{A}{R}{0{,}6}{0{,}4}{0{,}9}{0{,}1}{0{,}7}{0{,}3}
  \feuilles{A}{R}
\end{tikzpicture}
\end{minipage}
\cas{2)}
\begin{calculs}
P(A\cap R) &= P(A)\times P_A(R) &&\justif{chemin $A$ puis $R$}\\
P(A\cap R) &= 0{,}6\times 0{,}9 &&\\
P(A\cap R) &= \res{0{,}54} &&
\end{calculs}
\cas{3)}
\begin{calculs}
P(R) &= P(A\cap R)+P(\overline{A}\cap R) &&\justif{formule des probabilités totales}\\
P(R) &= 0{,}54+0{,}4\times 0{,}7 &&\justif{$P(A\cap R)$ calculé en 2)}\\
P(R) &= 0{,}54+0{,}28 &&\\
P(R) &= \res{0{,}82} &&
\end{calculs}
\cas{4)}
\explique{On note $E$ \og le voyageur utilise les deux moyens de transport\fg{} : train puis
bus ($A\cap\overline{R}$), ou bus puis train ($\overline{A}\cap R$).}{deux chemins de
l'arbre, qu'on ne peut pas suivre en même temps}
\begin{calculs}
P(E) &= P(A\cap\overline{R})+P(\overline{A}\cap R) &&\justif{on additionne les deux
chemins}\\
P(E) &= 0{,}6\times 0{,}1+0{,}4\times 0{,}7 &&\\
P(E) &= 0{,}06+0{,}28 &&\\
P(E) &= \res{0{,}34} &&
\end{calculs}
\cas{5)}
\begin{calculs}
P_R(A) &= \dfrac{P(R\cap A)}{P(R)} &&\justif{on sait que $R$ est réalisé : propriété 4}\\
P_R(A) &= \dfrac{0{,}54}{0{,}82} &&\justif{questions 2 et 3}\\
P_R(A) &\approx \res{0{,}66} &&\justif{$0{,}6585\dots$ arrondi au centième}
\end{calculs}
\end{exo}

\partie{Événements indépendants}

\begin{exo}[Indépendants ou pas ?]
\cas{a)}
\begin{calculs}
P_F(N) &= \dfrac{20}{100} = \res{0{,}2} &&\justif{colonne $F$ : $100$ vélos, dont $20$ avec
un pneu à changer}\\
P(N) &= \dfrac{50}{250} = \res{0{,}2} &&\justif{sur l'ensemble des $250$ vélos}
\end{calculs}
\explique{Conclusion : $P_F(N)=P(N)$, donc \res{$F$ et $N$ sont indépendants}.}{savoir que
les freins sont à régler ne change pas la probabilité d'avoir un pneu à changer}
\cas{b)}
\explique{$P_R(B)=0{,}3$ et $P(B)=0{,}2$.}{\og un jour de pluie\fg{} : sachant $R$ ; \og sur
l'année\fg{} : $P(B)$}
\explique{Conclusion : $P_R(B)\neq P(B)$, donc \res{$R$ et $B$ ne sont pas
indépendants} : la pluie augmente la probabilité de venir en bus.}{une différence suffit}
\cas{c)}
\explique{On note $F_1$ \og la première lampe fonctionne\fg{} et $F_2$ \og la seconde lampe
fonctionne\fg{}.}{on nomme les événements}
\begin{calculs}
P(F_1\cap F_2) &= P(F_1)\times P(F_2) &&\justif{$F_1$ et $F_2$ sont indépendants
(propriété 7)}\\
P(F_1\cap F_2) &= 0{,}98\times 0{,}9 &&\\
P(F_1\cap F_2) &= \res{0{,}882} &&\\
P(\overline{F_1\cap F_2}) &= 1-0{,}882 = \res{0{,}118} &&\justif{l'éclairage n'est pas
complet : contraire}
\end{calculs}
\explique{L'éclairage reste complet pendant un an avec une probabilité de $0{,}882$ ; il ne
le reste pas avec une probabilité de $0{,}118$.}{on traduit par une phrase}
\end{exo}

\begin{exo}[Une pièce et un dé]
\cas{1)}
\begin{minipage}[c]{0.58\linewidth}\raggedright
Pièce : $2$ branches ; au bout de chacune, dé : $6$ branches.
\begin{calculs}
2\times 6 &= \res{12} &&\justif{$2$ choix, puis $6$ choix}
\end{calculs}
Il y a $12$ issues.
\end{minipage}\hfill
\begin{minipage}[c]{0.38\linewidth}\centering
\begin{tikzpicture}[x=1cm,y=0.32cm,font=\small]
  \node (O) at (0,0) {};
  \foreach \ca/\ya in {P/3.5,F/-3.5}{%
    \node[evt] (N\ca) at (1.2,\ya) {\ca};
    \draw (O.center) -- (N\ca);
    \foreach \cb/\db in {1/2.5,2/1.5,3/0.5,4/-0.5,5/-1.5,6/-2.5}{%
      \node[evt,font=\footnotesize] (N\ca\cb) at (2.4,\ya+\db) {$\cb$};
      \draw (N\ca) -- (N\ca\cb);
      \node[right,font=\footnotesize] at (N\ca\cb.east) {\ca\cb};}}
\end{tikzpicture}
\end{minipage}
\cas{2)}
\explique{La pièce et le dé sont équilibrés : les $12$ issues sont équiprobables.}{chaque
chemin a la même chance}
\begin{calculs}
P(A) &= \dfrac{6}{12} = \res{\dfrac{1}{2}} &&\justif{les $6$ issues qui commencent par P}\\
P(B) &= \dfrac{2}{12} = \res{\dfrac{1}{6}} &&\justif{P6 et F6}\\
P(A\cap B) &= \res{\dfrac{1}{12}} &&\justif{une seule issue : P6}
\end{calculs}
\cas{3)}
\begin{calculs}
P_A(B) &= \dfrac{P(A\cap B)}{P(A)} &&\justif{propriété 4}\\
P_A(B) &= \dfrac{1}{12}\div\dfrac{1}{2} &&\\
P_A(B) &= \dfrac{1}{12}\times 2 &&\justif{diviser par $\frac{1}{2}$, c'est multiplier par
$2$}\\
P_A(B) &= \dfrac{2}{12} = \res{\dfrac{1}{6}} &&\justif{on simplifie par $2$}
\end{calculs}
\explique{Conclusion : $P_A(B)=P(B)$, donc \res{$A$ et $B$ sont indépendants}.}{le
résultat de la pièce ne change rien au dé}
\rem{On le voit aussi sur l'arbre : parmi les $6$ issues qui commencent par P, une seule
finit par $6$.}
\cas{4)}
\begin{calculs}
P(A)\times P(B) &= \dfrac{1}{2}\times\dfrac{1}{6} = \res{\dfrac{1}{12}} &&\justif{on
multiplie les numérateurs et les dénominateurs}
\end{calculs}
\explique{On retrouve $P(A\cap B)=\dfrac{1}{12}$ (question 2).}{propriété 7}
\end{exo}

\begin{exo}[L'imprimante 3D]
\cas{1)}
\begin{calculs}
P(A\cap B) &= P(A)\times P(B) &&\justif{$A$ et $B$ sont indépendants (propriété 7)}\\
P(A\cap B) &= 0{,}05\times 0{,}02 &&\\
P(A\cap B) &= \res{0{,}001} &&
\end{calculs}
\cas{2)}
\explique{L'impression s'arrête si $A$ ou $B$ est réalisé : on cherche $P(A\cup B)$.}{au
moins une des deux pannes}
\begin{calculs}
P(A\cup B) &= P(A)+P(B)-P(A\cap B) &&\justif{formule de la réunion}\\
P(A\cup B) &= 0{,}05+0{,}02-0{,}001 &&\\
P(A\cup B) &= \res{0{,}069} &&
\end{calculs}
\cas{3)}
\begin{calculs}
P(\overline{A\cup B}) &= 1-0{,}069 = \res{0{,}931} &&\justif{aller jusqu'au bout : contraire
de s'arrêter}
\end{calculs}
\explique{Une impression va jusqu'au bout avec une probabilité de $0{,}931$, soit
$93{,}1\,\%$.}{on traduit par une phrase}
\end{exo}

\end{document}
